\bookchapter{第二十七章：お父さんお母さんもラクラク——スマホのシニア向け完全ガイド}

\chapterintro{正直に言うと、この章を書いている時に両親のことを思い出した。母は“システム更新”のポップアップが出るたびに、スマホを壊すのが怖くてタップできない。父のWeChatの文字は蟻より小さいが、大きくする方法を知らず、二年間目を細めて見ていた。これは“バカ”だからではなく、スマホのデフォルトが全員を二十代だと思い込んでいるからだ。この章では、一番分かりやすい言葉で、お年寄りがスマホを使う際のもっとも頭痛の種を一つずつ解決していく。}

\booksection{まず心構えを正そう：これは“お年寄り使いこなせない”ではなく“スマホが作れていない”のだ}

設定をいじる前に、一つの考え方を明確にしておこう：\textbf{現在のスマホのUIは、設計段階からあまり高齢ユーザーを考慮していない}。ポップアップは小さくて見えない、操作ジェスチャーは正確なタッチを要する、アイコンがどれも似ている——これらはユーザーの問題ではなくプロダクトの問題だ。だからお年寄りにスマホを教える際、最初にすべきは教えることではなく、環境を整えることだ。\textbf{環境を整える}。環境さえ整えば、問題の大半は消えてしまう。

お年寄りにスマホを教える三つの鉄則：(1)\textbf{急かさない}——あなたは十年スマホを使っているが、彼らはまだ三ヶ月かもしれない。(2)\textbf{例えを使う}——“Wi-Fiは家の水道管のようなもの、蛇口をひねって初めて水が出る”、こういう例えは専門用語より一万倍分かりやすい。(3)\textbf{書き留める}——口頭で一度教えても忘れるかもしれない。紙に重要ステップを書いてスマホケースの裏に貼っておくのが一番効く。

\booksection{一、第一歩：「簡単モード」をオンにする——スマホがスッキリする}

現在のスマホにはほぼ“簡単モード”（機種によっては“ケアモード”“シニアモード”“お年寄りモード”）が内蔵されており、オンにするとスマホ全体が大きく変わる：

\begin{itemize}
  \item \textbf{文字とアイコンが大きくなる}——もう老眼鏡なしで見られる。
  \item \textbf{ホーム画面のアイコンが大幅に整理される}——電話、SMS、カメラ、アルバム、WeChatなど最もよく使うAppだけを残し、残りはすべて隠す。
  \item \textbf{操作がシンプルになる}——派手な機能（左端のパネル、スマートアシスタント、フローティングウィンドウなど）はデフォルトでオフになり、UIはスッキリ志向。
  \item \textbf{音が大きくなる}——着信音、通知音、通話音量がすべて自動で最大になる。
\end{itemize}

\booksubsection{各ブランドの簡単モードのオンにする方法（そのまま操作すればOK）：}

\begin{center}
\begin{tabularx}{\linewidth}{|X|X|X|}
\hline
\textbf{ブランド} & \textbf{モード名} & \textbf{オンにする方法} \tabularnewline
\hline
iPhone & アクセシビリティ（専用の簡単モードはないが、個別に大きく設定可能） & 設定 > 表示と明るさ > 文字サイズ、最大までドラッグ。さらに"アクセシビリティ > 表示と文字サイズ"で"太字テキスト""ボタン形状"をオン \tabularnewline
Apple急速充電 & 簡単モード & 設定 > システムと更新 > 簡単モード > オン \tabularnewline
Xiaomi/Redmi & シンプルモード & 設定 > その他の設定 > シンプルモード > オン \tabularnewline
OPPO/OnePlus & 簡単モード & 設定 > 特殊機能 > 簡単モード \tabularnewline
vivo/iQOO & シンプルデスクトップ & 設定 > デスクトップ、ロック画面と壁紙 > シンプルデスクトップ \tabularnewline
Samsung & 簡単モード & 設定 > 表示 > 簡単モード \tabularnewline
\hline
\end{tabularx}
\end{center}

他のブランドを使っている場合、“設定”上部の検索ボックスに“簡単”または“シニア”と入力すればだいたい見つかる。見つからない場合は、下記の“手動で大きくする”方法を一つずつ設定すれば同じ効果が得られる。

\booksection{二、簡単モードを使わなくても手動でOK——“大きな文字+大きなアイコン”の五ステップ}

お年寄りによっては簡単モードのUIを“年寄りくさい”と嫌がる。その場合は以下の五項目を手動で調整すれば同じ効果がある：

\begin{enumerate}
  \item \textbf{文字を最大にする}：設定 > 表示 > 文字サイズ、一番右までドラッグ。ついでに"表示サイズ"（機種により"表示スケール"）も最大にすると、すべてのUI要素が一緒に大きくなる。
  \item \textbf{アイコンを大きくする}：ホーム画面で二本指でピンチ（写真を縮小するように）してデスクトップ設定 > デスクトップレイアウト（またはアイコンサイズ）に入り、マス目が最も少ないレイアウト（例：3×5）を選ぶと、アイコンが自然に大きくなる。
  \item \textbf{"高コントラストテキスト"をオンにする}：文字の輪郭をよりシャープにする。"アクセシビリティ"の中にこのスイッチがある。
  \item \textbf{不要な通知をオフにする}：設定 > 通知、電話、SMS、WeChat以外のApp通知をすべてオフにする。これでスマホがいちいち鳴らず、うるさくなくなる。
  \item \textbf{よく使う連絡先をホーム画面に置く}：ホーム画面の空白部分を長押し > ウィジェット追加 > 連絡先、配偶者、子供、医者、緊急連絡先のクイックダイヤルをホーム画面に置く。ワンタップで発信でき、連絡帳を探す必要がない。
\end{enumerate}

\begin{tcbpractice}{実践：お年寄りのスマホを"リフォーム"しよう}
\begin{enumerate}
  \item 簡単モードをオンにする（上記のブランド別表通りに操作）。
  \item バッテリー残量の数値表示をオン（設定 > バッテリー > バッテリー残量％）、残量がいつでも分かるようにする。
  \item 通知パネルを開き、Bluetooth、NFC、機内モードなどお年寄りには不要なクイックスイッチを削除または隠す。
  \item ホーム画面に4つのクイックダイヤルを置く：配偶者、子供、かかりつけ医（いれば）、社区/管理事務所。
\end{enumerate}

\end{tcbpractice}

\booksection{三、詐欺防止——この章は必ずお年寄りと一緒に読もう}

高齢者はテレコミ詐欺の最大の被害層だ。警戒心が低いのではなく、詐欺師の手口が綿密に設計され、高齢者の心理的弱点を狙っているからだ。以下は私が整理した最もよくある詐欺の手口と対処法だ。お年寄りに直接読み聞かせるか、印刷して電話の横に貼っておくことをお勧めする。

\booksubsection{六大よくある詐欺（見たらすぐ電話を切る）：}

\begin{itemize}
  \item \textbf{(1) "お孫さんが事故に遭った、すぐ送金して"}——詐欺師は"お孫さんが外でトラブルを起こして警察に捕まった、保証金が必要"と言う。覚えておこう：\textbf{送金しないと"釈放しない"という電話はすべて詐欺だ}。公安·検察·裁判所が電話で金銭を要求することはない。
  \item \textbf{(2) "あなたの医療保険カードが停止されました"}——SMSで医療保険カードに異常があると言い、リンクをタップして"有効化"させる。覚えておこう：\textbf{医療保険局、社会保障局がSMSでリンクを促すことはない}。この場合は12333（社保公式ホットライン）に電話して確認する。
  \item \textbf{(3) "高収益の財テク、絶対に儲かる"}——WeChatグループやSMSで各種"シニア財テク""国家支援プロジェクト"を推奨し、高額リターンを約束する。覚えておこう：\textbf{年利5\%以上で"元本保証"を謳うものは99.9\%詐欺だ}。
  \item \textbf{(4) "私はあなたの上司/古い友人だ、番号が変わった"}——知人を装って金を借りる。覚えておこう：\textbf{電話して確認するのが一番効く}、"電話に出られない"と言うやつは八割詐欺師だ。
  \item \textbf{(5) "無料で卵/健康食品の講座"}——線下の無料イベントで"健康講座"を聴かせ、法外な価格の健康食品を売りつける。覚えておこう：\textbf{棚から牡丹餅は落ちない、無料イベントの裏には商売がある}。
  \item \textbf{(6) "スマホがウイルス感染した、リモート修復が必要"}——詐欺師がリモートコントロールApp（向日葵、TeamViewerなど）をダウンロードさせ、リモートでスマホを操作して金を奪う。覚えておこう：\textbf{誰にもリモートでスマホを操作させてはいけない！}スマホメーカー、通信キャリア、銀行が電話でリモートコントロールソフトのインストールを促すことはない。
\end{itemize}

\booksubsection{三重の防御ラインを敷けば、90\%の詐欺を防げる：}

\begin{enumerate}
  \item \textbf{第一のライン：スマホ内蔵のブロック機能}——Androidスマホは"迷惑ブロック"をオン（電話Appまたはスマホ管理App内）、iPhoneは"不明な発信者を消音"をオン（設定 > 電話 > 不明な発信者を消音）。
  \item \textbf{第二のライン：キャリアのブロック}——SMSで"KTFSR"を10086に送信、China Mobileの高頻度迷惑電話ブロックを無料で有効化。China TelecomやChina Unicomにも類似サービスあり。キャリアのカスタマーサービスに電話して開通してもらえばよい。
  \item \textbf{第三のライン：国家反詐欺センターApp}——アプリストアで"国家反詐騙センター"Appをダウンロード、登録して着信警告をオンにする。このAppは公安のデータベースと連携し、詐欺師の番号がかかってくると即座に識別する。
\end{enumerate}

\begin{tcbwarn}{お年寄り向け"防詐欺のお呪い"（印刷して電話の横に貼る）}
\begin{itemize}
  \item \textbf{知らない電話で送金要求 切る！}
  \item \textbf{SMSのリンクをタップさせようとする 削除！}
  \item \textbf{高収益·元本保証と言う ブロック！}
  \item \textbf{リモートコントロールソフトを入れさせようとする 即切る！}
  \item \textbf{本物の警察は電話で事件を処理せず、ましてや金を要求しない！}
  \item \textbf{迷ったら、まず子供に電話して確認しよう。}
\end{itemize}

\end{tcbwarn}

\booksection{四、リモートサポート——離れていても親のスマホを修理できる}

お年寄りが問題に遭遇した時、電話でいくら説明しても該当するボタンが見つからないことがある。そんな時に“リモートサポート”機能が役立つ——自分のスマホから相手のスマホ画面を見て、操作を手伝える。各ブランドの方法は異なる：

\begin{itemize}
  \item \textbf{Huawei/Honor}：内蔵の"Changlian"アプリに"画面共有"機能がある。親子でChanglianを入れ、ビデオ通話を開始して"画面共有"をタップすれば、相手の画面に丸を描いて"ここをタップ"と教えられる。
  \item \textbf{Xiaomi/Redmi}：Xiaomiのビデオ通話Appに"リモートサポート"があり、画面共有+音声通話の同時進行に対応。
  \item \textbf{：一部機種はアンロック可能、申請が必要。}：内蔵の"ファミリースペース"または"リモートガード"、設定で"リモート"と検索すれば見つかる。
  \item \textbf{汎用ソリューション（全ブランド対応）}："向日葵リモートコントロール"Appをダウンロード（無料版で十分）、双方にインストール。同じアカウントで登録するか友達追加し、親にスマホに表示された認証コードを読んでもらえば、リモートで画面を見て操作できる。
  \item \textbf{：iOSメモ標準のスキャン機能がすでに十分優れています。メモを開く カメラアイコンをタップ 「書類をスキャン」を選択、書類に向けると自動的に枠を認識して透視を補正します。撮影後に色調整ができます——「白黒」モードを選ぶとスキャナーに最も近い効果になり、「グレースケール」は画像入りの資料に適しています。}：FaceTime通話中に"画面共有"に対応。FaceTime通話中に画面共有ボタンをタップすれば、相手があなたの画面を見られる。ただしiOSのセキュリティ制限によりリモート操作はできず、"指示型の指導"のみ。
\end{itemize}

\booksection{五、WeChatの“シニア改造”——お年寄りにこれだけ伝えればOK}

WeChatはお年寄りが最もよく使うAppだが、UIが密集しすぎている。以下は即効性のある調整：

\begin{itemize}
  \item \textbf{WeChat内の文字を大きくする}：WeChat > 私 > 設定 > 一般 > 文字サイズ、最大までドラッグ。これはWeChat内の文字にのみ影響し、システム全体の文字サイズには影響しない。
  \item \textbf{"ケアモード"をオンにする}：WeChat > 私 > 設定 > ケアモード > オン。オンにすると文字とボタンが大きくなり、色のコントラストが強化され、ワンタップで"文字メッセージを読み上げ"（テキストメッセージをタップするだけで読み上げられる）。
  \item \textbf{重要な連絡先に分かりやすいメモを付ける}：子供や親戚の名前を連絡先で分かりやすく変更する——例えば"長男-小明""次女-小紅"など、WeChatを開けば誰だかすぐ分かるように。
  \item \textbf{"音声入力"を教える（タイピングの代わり）}：チャット欄のキーボードで、スペースキーを長押し（Android）またはマイクアイコンをタップ（iPhone）、直接話しかければ文字になって送信される。認識精度はすでに非常に高く、ピンイン入力ができないお年寄りには"画期的"な便利さだ。
  \item \textbf{"発見"ページの不要な機能をオフにする}：WeChat > 私 > 設定 > 一般 > 発見ページ管理、"モーメンツ""動画アカウント""ライブ""ショッピング""ゲーム"など使わないものをすべてオフにし、干渉を減らす。
\end{itemize}

\booksection{六、スマホを失くしたら？お年寄り向け“緊急対応プラン”を作る}

備えあれば憂いなし。スマホを失くす前に以下を済ませておこう：

\begin{enumerate}
  \item \textbf{"スマホを探す"がオンになっているか確認}：iPhoneは設定 > あなたの名前 > 探す > iPhoneを探す；Androidは設定で"デバイスを探す"を検索。このスイッチがオンになっていることを必ず確認。
  \item \textbf{重要なパスワードとアカウントを紙に書く}：Apple ID / Huaweiアカウント / Xiaomiアカウントのパスワードを家の引き出しに一部、信頼できる子供に一部。スマホ紛失後にこれらのアカウントにログインしてリモートロックと位置追跡が必要になるので、パスワードを忘れると大変だ。
  \item \textbf{お年寄りに一言教える}：\textbf{"スマホを失くしたら、すぐ家族に電話して。"}家族はすぐにiCloudまたはAndroidのデバイス探索Webページにログインし、リモートで"紛失モード"を設定できる。そうすれば拾った人も開けず、使えない。
\end{enumerate}

\booksection{七、ワイヤレスイヤホンと補聴器の互換性——見落とされがちな便利機能}

多くのお年寄りは聴力に問題があるが、スマホが補聴器と連携できることを知らない。iPhoneと一部のAndroidスマホはMFi（Made for iPhone）補聴器に対応し、通話音声を直接補聴器に送れる。スマホの受話器を耳に当てるよりずっとクリアだ。

\begin{itemize}
  \item \textbf{接続方法}：iPhoneの設定 > アクセシビリティ > 補聴器、Bluetoothをオンにし、補聴器のペアリング手順に従う。
  \item \textbf{補聴器を使わない場合}："アクセシビリティ > オーディオ/視覚"で"通話ノイズキャンセル"と"モノラルオーディオ"をオンにすると、通話の明瞭度が改善する。
  \item \textbf{Android}："アクセシビリティ > 聴力強化"に類似の設定がある。ブランドにより位置が少し異なるので"補聴"で検索する。
\end{itemize}

\booksection{八、健康関連機能——スマホは通信ツールにとどまらない}

\begin{itemize}
  \item \textbf{SOS緊急求助}：電源ボタンを5回連続押し（大部分のスマホ）、自動的に緊急電話を発信し、設定済みの緊急連絡先に位置情報を送信する。お年寄りのために事前に設定しておこう——iPhone："設定 > 緊急SOS"、Android："設定 > セキュリティ > 緊急求助"。教える前に一度リハーサルし、本番で焦って押し間違えないようにする。
  \item \textbf{転倒検出}：Apple Watchにこの機能があり、転倒を検出してしばらく動かないと、自動的に救急電話を発信し家族に通知する。予算があればお年寄りに一台贈ろう。緊急の瞬間に命を救える。
  \item \textbf{服薬リマインダー}：iPhoneの"ヘルスケア"Appに服薬リマインダー機能が内蔵。薬名、時間、服用量を設定すれば、時間になるとリマインドが表示される。Androidは"繰り返しアラーム"や"服薬リマインダー"系Appをダウンロードする。
\end{itemize}

\begin{tcbsummary}{本章のまとめ}
\begin{itemize}
  \item まず簡単モードをオンにするか手動で文字/アイコンを大きくし、環境を整えれば問題の大半は解決する
  \item お年寄りにスマホを教える三鉄則：急かさない、例えを使う、書き留める
  \item 六大よくある詐欺は繰り返し注意し、三重の防御ラインで90\%の詐欺を防げる
  \item リモートサポートで離れていても親のスマホを修理できる
  \item WeChatの"ケアモード"+"音声入力"がお年寄りに最も実用的な二つの機能
  \item SOS緊急求助と転倒検出は、緊急の瞬間に命を救える設定。5分かけて設定する価値がある
\end{itemize}

\end{tcbsummary}
